% !PRINTED 588
\weberpage{588}
%<<CONT
plex multiplication of the elliptic functions gives examples, the
partial fundamental ideal $\mathfrak{G}_R$ is to be set equal to 1.

If the $n$ fields (3) are identical with one another, then $\Omega$ is called a
\gesperrt{normal field with respect to} $R$ (relatively normal, if $R$ is not the absolute
domain of rationality). The totality of the $n$ substitutions $(\Theta, \Theta_n)$
that effect the transition of the numbers of the field $\Omega$ to the conjugate
numbers forms a group of degree $n$ --- as we have already seen in \S.\,147 of the
first volume --- which is nothing other than the \gesperrt{Galois} group of the
equation (1) in the domain of rationality $R$, and which we here call the \gesperrt{group of
the field $\Omega$ (with respect to $R$)}.

If in particular this group is commutative, then $\Omega$ is an
\gesperrt{Abelian field with respect to} $R$ (a relatively \gesperrt{Abelian field}). In this
case (1) is an \gesperrt{Abelian} equation in the domain of rationality $R$.

\begin{center}
\S.\,160.\\[6pt]
\gesperrt{Prime ideals in the relatively normal field.}
\end{center}

We now assume that the field $\Omega$ is normal with respect to the arbitrary
field $R$, and under this assumption we investigate the prime ideals of the field
$\Omega$.

Let the group of $\Omega$ with respect to $R$ be $\Phi$, and let the substitutions
of $\Phi$ be denoted by $\varphi, \varphi_1, \ldots$ We further denote by
\weq{1}{
\omega \mid \varphi
}
\gesperrt{the number that arises from $\omega$ by the substitution $\varphi$}, so that $\omega$
and $\omega \mid \varphi$ are contained in $\Omega$.

Just as the numbers $\omega$ do, so also all functionals, and with them at the same
time all ideals $\mathfrak{a}$ of the field $\Omega$, pass by a substitution
$\varphi$ into definite functionals and ideals $\mathfrak{a} \mid \varphi$; and if
$\omega$ is an integer divisible by $\mathfrak{a}$, then $\omega \mid \varphi$ is
divisible by $\mathfrak{a} \mid \varphi$.

If $\mathfrak{a}$ is any ideal in $\Omega$, then there are certain substitutions
$\psi$ in $\Phi$ that satisfy the condition
\weq{2}{
\mathfrak{a} \mid \psi = \mathfrak{a}
}
among them always the identical substitution, and these substitutions $\psi$
form a group $\Psi$, which is a divisor of $\Phi$
%CONT>>
